% Act 3, rule one: what a reused state must hold depends on the next step. One table, one status line.
\begin{frame}[t]{\ev{\tagM\,\tagP}\surf{COMPUTE ONCE}The next step decides which state to reuse.}
\vspace{.1cm}
{\fontsize{14}{17}\selectfont\renewcommand{\arraystretch}{1.32}
\begin{tabularx}{\linewidth}{@{}L{5.1cm}L{3.6cm}Y@{}}
\toprule
\textcolor{Teal}{\textbf{Next step}} & \textcolor{Teal}{\textbf{Reuse by}} & \textcolor{Teal}{\textbf{Supplies}}\\
\midrule
Rebuild & \textbf{Cache} & files\\
Continue a live run & \textbf{Restore} & its memory\\
Reuse an answer & \textbf{Replay} & the recording\\
\bottomrule
\end{tabularx}\par}
\vspace{.4cm}
{\fontsize{13}{17}\selectfont\tagM\enspace Restore fan-out on islo: \textbf{255/256} succeeded.\par
\vspace{.12cm}\tagP\enspace Live memory: untested (Gate~0).\par}
\sources{Restore fan-out: many restores of one named snapshot (\href{https://arxiv.org/abs/2607.09689}{arXiv:2607.09689v4}, Table~2). See also CRIU, Firecracker, rr.}
\qadetail{Fork semantics are backend-specific: filesystem clone, artifact cache, process fork, guest-memory restore and remote clone differ. A fork supplies several continuations from one captured parent with private writable state; fork identity is distinct from checkpoint identity and observation identity, and authority stays outside child-owned state. The three rows are the three ways to supply what the next step reads: rebuild from a cache (worktrees parallelize edits without preserving live processes), restore a checkpoint, or replay a record. Choose the cheapest mechanism that supplies what the next step reads. The next slide prices a restore against a warm cache for a next step that reads only files, so a restore is compared with a warm cache, not with a cold start; a warm cache cannot continue a live run. Process checkpoints (CRIU, gVisor save/restore) and VM snapshots (Xen, SnowFlock, Firecracker) preserve only their declared state surfaces; the slide groups them without equating their supported state. Replay (rr, DSec) reuses a recorded response without reexecuting a remote service. Restore fan-out, measured: a named 141 MB islo snapshot, restore--run--capture round trips, 255 of 256 succeeded at concurrency 12, p50 6.87 s, p95 9.04 s, client-measured, teardown excluded from the percentiles (arXiv:2607.09689v4, Table 2). These are API round trips, not mechanism latency or a provider ranking, and they are separate from the factory runs, which used Docker sandboxes, serially, with nothing forked. A snapshot name does not establish memory semantics: memory preservation on this path is unverified, so the supported description is restore fan-out, not a memory fork; the Gate 0 RAM-only 128-bit nonce probe is preregistered and unrun. VM and process checkpoints can capture memory by design; the open question is only whether the path we restore keeps it. Network, time, randomness and remote model sessions need declared semantics; hosted-model weights, KV state and sampling stay outside the guest. The factory runs established files, dependencies and failed-test feedback, not a necessary live process; a running application at a reproduced failure is a hypothetical future workload. Editing loaded code may require reload or restart and invalidate the captured live state. SnowFlock: Lagar-Cavilla et al., EuroSys 2009, built on Xen.}
\note{[1:05] To compute once, first ask what the reused state must hold. That depends on the next step. If it edits code and rebuilds, the cache we saw earlier supplies the files. If it continues a live run, we restore a process or VM checkpoint and get only the memory that mechanism keeps. If it needs an earlier answer, replay returns the recording. What runs today is restore fan-out on islo, our sandbox service: many restores of one named snapshot. Two hundred fifty-five of two hundred fifty-six restore-run-capture calls succeeded. Whether live memory survives is untested until Gate zero, our first preregistered probe. So pick the cheapest mechanism that supplies what the next step reads. For a step that only reads files, a restore and a warm cache compete.}
\end{frame}
